% Minimal preamble for P2 cards taken from sources/comparability-partitions-2408.13523v2/source0.tex
% (arXiv:2408.13523v2, "Partitioning perfect graphs into comparability graphs").
% Every line below is copied mechanically from that file; nothing was invented, renamed or redefined.
%   line 2    : \documentclass[12pt]{article}
%   line 3    : \usepackage{amssymb,latexsym}
%   line 4    : \usepackage{amsmath}
%   line 5    : \usepackage[dvipdfm]{graphicx}  -- the dvipdfm driver option is dropped, because the
%               assembler forces the xelatex engine and a forced driver option contradicts it.
%   line 16   : \usepackage{amsfonts}
%   line 17   : \usepackage{amsthm}
%   lines 20-30: the author's theorem declarations
%   lines 32-41: the author's own macro definitions
% Omitted source preamble lines and why they are not needed for these body spans:
%   line 1 (provenance comment), lines 6-11 (tikz and its libraries: no drawing in the retained spans),
%   lines 7 and 9 (verbatim), lines 12-13 (page geometry), line 14 (rawfonts),
%   line 15 (psfrag), line 18 (hyperref with the pagebackref option: the retained spans print no
%   citation anchors and the assembly adds its own hyperref machinery), lines 43-61 (the author's
%   backref neutralisation and his margin-note macros, which depend on the omitted hyperref).
% The retained spans of the body use only \em, \bf, \cal, amsthm environments, amsmath displays and
% the macros defined in lines 32-41 below.
\documentclass[12pt]{article}
\usepackage{amssymb,latexsym}
\usepackage{amsmath}
\usepackage{graphicx}
\usepackage{amsfonts}
\usepackage{amsthm}

\newtheorem{theorem}{Theorem}%[section]
\newtheorem{problem}[theorem]{Problem}%[section]
\newtheorem{observation}[theorem]{Observation}
\newtheorem{lemma}[theorem]{Lemma}
\newtheorem{corollary}[theorem]{Corollary}
\newtheorem{definition}{Definition}
\newtheorem{example}[theorem]{Example}
\newtheorem{conjecture}{Conjecture}%[section]
\newtheorem{proposition}[theorem]{Proposition}
\newtheorem{claim}[theorem]{Claim}
\newtheorem{remark}[theorem]{Remark}

\def\qed{\ifhmode\unskip\nobreak\fi\quad\ifmmode\Box\else$\Box$\fi}
\def\pont{\hspace{-6pt}{\bf.\ }}
\def\qed{\hfill $\Box$}

\def\R{{\mbox{\boldmath$\mathcal R$}}}
% the author's second definition of \R, left commented out exactly as in the source
\def\F{{\mathcal F}}
\def\P{{\mathcal P}}

\def\newpage{\vfill\eject}
\title{Partitioning the disjointness graph of the subtrees of a tree into two comparability graphs}
\author{András Gyárfás, Márton Marits, Géza Tóth}
\date{arxiv:2408.13523v2; CC BY 4.0}
\begin{document}
\maketitle
\noindent\textit{Editorial note.} The counted unit is the article's own Theorem co-triangulated at source lines 313-316 together with the author's proof that opens with the bold header Proof. at source line 319 and closes with the square at source line 373; the header is the article's own, the pairing is by adjacency through the blank line and the medskip at source lines 317-318, and the proof ends before the figure environment that begins at source line 377.\par\medskip
\noindent\textit{Editorial note.} The retained context at source lines 117-120 and 128-129 is the article's own description of a comparability graph and its own definitions of the parameters p(G) and c(G); the retained context at source lines 276, 278 and 309-311 is its own description of subtree intersection graphs, of chordal alias triangulated graphs and of the complements of interval graphs.\par\medskip
\noindent\textit{Editorial note.} The source preamble cannot be retained as it stands: it loads rawfonts, psfrag, verbatim, tikz with several libraries, hyperref with the pagebackref option and a dvi driver option for graphicx, none of which the retained body spans use, and it redefines backref commands that only exist because of that hyperref option. The retained preamble file therefore carries, line for line, source line 2, source line 3, source line 4, source line 16, source line 17, the theorem declarations of source lines 20-30 and the author macro definitions of source lines 32-41; the only textual change anywhere is the removal of the dvipdfm driver option from the graphicx line, because the assembly forces the xelatex engine. No new macro, package or declaration was introduced.\par\medskip
\noindent\textit{Editorial note.} No retained line contains a citation: the sentences of the article that cite Dilworth, Mirsky, Gavril, Walter and the surveys are deliberately not retained, so the material carries no bibliography.\par\medskip
\noindent\textit{Editorial note.} The article's theorems gensplit, bipline, unit, loglog and ub, its lemma rec and its proposition H9 are not counted here and are not used as premises.\par\medskip
\noindent\textit{Allowed standard prerequisites (editorial).}\begin{itemize}
\item Comparability graph of a partially ordered set. Given a partially ordered set, its comparability graph has the elements as vertices and an undirected edge between two vertices exactly when the corresponding elements are comparable; the article's own description is retained at source lines 117-120.. Fixes the object that the proof constructs twice on the disjoint pairs of a subtree system, and whose union must reproduce the disjointness graph of the statement.
\item The partition number p(G) and the cover number c(G). p(G) is the minimum m such that the edge set of G is the union of m edge-disjoint comparability subgraphs of G, and c(G) is the same number without the edge-disjointness requirement, so that p(G) is at least c(G); retained at source lines 128-129.. The counted theorem asserts c(G) <= p(G) <= 2, so both parameters and their mutual relation must be fixed before the statement can be read.
\item Disjointness graph of a system of subtrees. For a tree T and a system F of subtrees of T, the disjointness graph has F as its vertex set and an edge between two subtrees exactly when they are disjoint; a graph is co-triangulated when it arises in this way, and the article records at source lines 276 and 278 that the complements of these graphs are the chordal, alias triangulated, graphs.. Identifies the graph G of the counted theorem with the disjointness graph of F, which is the graph the two constructed partial orders must partition.
\end{itemize}\medskip
Here we focus on a well known family of perfect graphs, the {\em comparability graphs},
which can be obtained from a partially ordered set as follows.
Vertices correspond to the elements of the partially ordered set,
and two vertices are connected by an undirected edge if and only if the corresponding elements

Let $p(G)$ denote the minimum number $m$ such that $E(G)=\cup_{i=1}^m E(G_i)$ where the $G_i$-s are edge-disjoint comparability subgraphs of $G$.
Dropping  the condition that the $G_i$-s are edge-disjoint, we define $c(G)$ as the smallest number of comparability subgraphs of $G$ needed to {\em cover} $E(G)$.

A more general family of perfect graphs is the intersection graphs of subtrees of a tree.

not contain induced cycles of length at least four and called {\em chordal, alias triangulated graphs}.

The complements of interval graphs, i.e. the disjointness graphs of interval systems are comparability graphs.
This is not true for its superclass, the complements  of the subtree graphs.
Figure 1 shows a perfect graph which is not a comparability graph but can be represented

\begin{theorem} \label{co-triangulated}
  Assume that $G$ is a co-triangulated graph,
  i.e. the disjoint\-ness graph of subtrees of a tree. Then $c(G)\le p(G)\le 2$.
\end{theorem}

\noindent{\bf Proof. } Let $T$ be a tree and $\cal{F}$ is a system of subtrees of $T$.
We construct two comparability graphs on the disjoint pairs of $\cal{F}$.
 Let $r$ be any vertex of $T$, the root. For any subtree $A\in {\cal{F}}$
there is a unique vertex of $A$ that is closest to $r$, call this vertex the root
of $A$.
For the first partial order $\prec_1$,
let $A,B\in {\cal{F}}$
be two disjoint subtrees. We say that
$A\prec_1 B$ if the unique path from the root of $B$ to $r$ contains the root of $A$. Intuitively,
we can say that $A\prec_1 B$ ``if $A$ is closer to $r$ than $B$''.
It is clear that $\prec_1$ is antisymmetric. Suppose that  $A\prec_1 B$ and  $B\prec_1 C$ for subtrees $A$, $B$, $C$.
Then the path $P_C$ from the root of $C$ to $r$ contains the root of $B$. But then $P_C$ contains the path $P_B$
from the root of $B$ to $r$. Since $P_B$ contains the root of $A$, $P_C$ also contains
the root of $A$. Moreover, the root of $B$ separates $A$ and $C$, so they are disjoint.
Therefore, $A\prec_1 C$, consequently, $\prec_1$ is transitive, so it is
a partial ordering on the disjoint pairs.


Let $N^{\le N}=\cup_{i=0}^NN^i$. We define a partial order $\prec_{lex}$ on $N^{\le N}$ as follows.
For $\alpha, \beta\in N^{\le N}$, let $\alpha\prec_{lex}\beta$ if none of $\alpha$ and $\beta$ is a prefix of the other,  and
$\alpha\prec\beta$ in the lexicographic ordering.

For the second partial order, we first define the labeling function $\psi : V(T) \rightarrow N^{\le N}$
in the following recursive way. In the root, we define $\psi= \emptyset$. For neighbors
of the root, we arbitrarily choose an ordering, and assign the numbers
$\{0,\ldots, k\}$
where $k$ is the degree of $r$ minus one.  In general, if a vertex $v$ has the label $\psi(v)$
and $v$ has at least one child, they get labels $\psi(v)0, \dots \psi(v)k$ where $k$ is the number of children minus one.

The value of $\psi$ at each
vertex is a sequence of natural numbers.
Observe that for disjoint subtrees $A,B$ with roots $a,b$, $A \prec_1 B$ if and only if
$\psi(a)$ is a prefix of $\psi(b)$.

Now let $A,B$ be two disjoint subtrees of $T$, with roots $a,b$ respectively,
which are not comparable by $\prec_1$.
Then let $A \prec_2 B$ if and only if  $\psi(a)\prec_{lex} \psi(b)$.


It is easily seen that $\prec_2$ is antisymmetric, it remains to prove
transitivity. Let $A,B,C$ be subtrees of $T$ with roots
$a,b,c$. If $A\prec_2 B$ and
$B \prec_2 C$, then clearly, $\psi(a) \prec_{lex}\psi(b)\prec_{lex}\psi(c)$.
This  implies that none of $\psi(a)$ and $\psi(c)$ is a prefix of the other and
$\psi(a) \prec_{lex} \psi(c)$.

Moreover, for any vertex $u$ of $A$, $\psi(a)$ is a prefix of $\psi(u)$,
therefore, $\psi(u)\prec_{lex}\psi(b)$. On the other hand,
for any vertex $v$ of $C$, $\psi(c)$
is a prefix of $\psi(v)$, therefore, $\psi(b)\prec_{lex}\psi(v)$.
Consequently,
$A$ and $C$
must be disjoint.
Therefore, the comparability graphs of $\prec_1$ and $\prec_2$ form
a partition of the disjointness graph of the subtrees of $T$.   \qed

\end{document}
